\chapter{Muscle Cramps}

\section*{Definition}
Sudden involuntary sustained contraction of muscle causing severe pain temporary incapacity; soreness persists after cramp ceases more likely tired muscles at night.

\section*{What Happens in Your Body}
Acute benign self-limiting differentiates from sustained spasm or tetany; dehydration electrolyte imbalance overuse pregnancy medication side effects common triggers; nocturnal cramps very common in elderly.

\section*{Causes}
\begin{itemize}
  \item Dehydration electrolyte loss sweating sports activity
  \item Muscle overuse exercise fatigue prolonged positioning
  \item Pregnancy hormonal changes fluid retention
  \item Medication side effects statins diuretics
  \item Peripheral artery disease circulation issues
\end{itemize}

\section*{When to See a Doctor}
See your doctor if:
\begin{itemize}
  \item Symptoms are severe or getting worse over time
  \item Cramps are persistent, frequent, or occur without an obvious cause
  \item You have other concerning symptoms such as fever, weight loss, or bleeding
\end{itemize}

\section*{Related Symptoms}
\begin{itemize}
  \item Muscle soreness after episode
  \item Localized swelling tenderness
  \item Visible muscle knotting
  \item Difficulty walking immediately after
  \item Recurrent pattern in same muscle
\end{itemize}
\textbf{Important:} If this symptom persists, your doctor will check for serious underlying causes.


\section*{Self-Care / Home Management}
\begin{itemize}
  \item Apply the RICE protocol: Rest, Ice (15-20 min every 2-3 hours), Compression with elastic bandage, Elevation above heart level.
  \item Gradual return to activity with gentle stretching once acute pain subsides.
  \item Over-the-counter anti-inflammatory medications (ibuprofen, naproxen) can reduce pain and swelling.
  \item Avoid activities that worsen symptoms until healing is well underway.
  \item Seek evaluation for severe pain, inability to bear weight, joint deformity, or symptoms lasting beyond 2 weeks.
\end{itemize}

\section*{How Your Doctor Will Evaluate You}
\begin{itemize}
  \item History focuses on mechanism of injury, onset (acute vs insidious), location, aggravating/relieving factors, and prior episodes.
  \item Physical examination includes inspection (swelling, bruising, deformity), palpation, range-of-motion, and strength testing.
  \item Imaging (X-ray for fracture, ultrasound for tendon/ligament injury, MRI for soft tissue detail) guides diagnosis.
  \item Functional assessment evaluates ability to perform daily activities and sports-specific movements.
\end{itemize}

\section*{Treatment Options}
\begin{itemize}
  \item Acute injuries benefit from RICE, activity modification, and gradual rehabilitation exercises.
  \item Physical therapy is central to recovery, focusing on range-of-motion, strengthening, and proprioceptive training.
  \item Corticosteroid injections may be considered for significant inflammatory conditions (tendinitis, bursitis).
  \item Orthopedic referral is indicated for complete tears, fractures, joint instability, or failure of conservative therapy for 6-8 weeks.
\end{itemize}

\clearpage
